\chapter{Related Work}
\label{ch:related}

This project draws on four lines of work: review-based recommendation and reviewed-item
retrieval, lexical and dense retrieval, explainable recommendation, and Thai language
processing and retrieval. Figure~\ref{fig:related-map} shows how they meet in this project, and
Table~\ref{tab:related-compare} compares the closest work with ours.

\begin{figure}[H]
\centering
\begin{tikzpicture}[font=\small,
  src/.style={draw,rounded corners,align=left,text width=4.7cm,inner sep=4pt,fill=gray!7},
  core/.style={draw,very thick,rounded corners,align=center,text width=2.9cm,inner sep=5pt,fill=blue!8},
  arrow/.style={-{Latex[length=2mm]},thick}]
\node[core] (c) {\textbf{This project}\\fair unit and volume\\control, evidence\\evaluation, Thai\\taste test set};
\node[src,above left=4mm and 5mm of c,anchor=south east] (a)
  {\textbf{Reviewed-item retrieval}\\RIRD \citep{rird2023}; multi-aspect fusion \citep{korikov2024multiaspect}; voting and expert search \citep{balog2006formal,macdonald2008voting}};
\node[src,below left=4mm and 5mm of c,anchor=north east] (b)
  {\textbf{Lexical and dense retrieval}\\BM25 \citep{robertson2009bm25}; BERT, SBERT, DPR \citep{devlin2019bert,reimers2019sbert,karpukhin2020dpr}; multilingual encoders \citep{reimers2020multilingual}; granularity \citep{chen2024densex}};
\node[src,above right=4mm and 5mm of c,anchor=south west] (d)
  {\textbf{Explainable recommendation}\\survey \citep{zhang2020explainable}; review sentences as explanations \citep{li2021extra}; faithfulness \citep{jacovi2020faithfulness,deyoung2020eraser}; user-centred quality \citep{zhang2025enhancing}};
\node[src,below right=4mm and 5mm of c,anchor=north west] (t)
  {\textbf{Thai NLP and retrieval}\\PyThaiNLP \citep{phatthiyaphaibun2023pythainlp}; WangchanBERTa \citep{lowphansirikul2021wangchanberta}; MIRACL \citep{zhang2023miracl}; negation in IR \citep{weller2024nevir}};
\draw[arrow] (a.south east) -- (c);
\draw[arrow] (b.north east) -- (c);
\draw[arrow] (d.south west) -- (c);
\draw[arrow] (t.north west) -- (c);
\end{tikzpicture}
\caption{The four lines of related work that this project combines.}
\label{fig:related-map}
\end{figure}

\section{Review-based recommendation}
Classic recommender systems model user--item interactions, for example with matrix
factorisation \citep{koren2009matrix}. These models predict what a user will like from past
behaviour, but they cannot answer a new free-text request and do not explain their output.
Review-based recommender systems add review text to the model. The survey by
\citet{hasan2025reviewsurvey} shows that most of this work uses reviews to improve rating
prediction or to build user and item profiles, and it names explainability and evaluation as
open problems. Our project uses reviews differently: as the searchable description of each item,
and as the evidence shown to the user.

Conversational recommender systems take the user's request into account by asking questions
and narrowing the set of candidates \citep{gao2021conversational}. When a request is vague,
generating questions about how the item will be used can help \citep{kostric2024usage}. Our
system is a single-query version of this idea. It accepts one natural-language request and, in
the live demo only, rewrites vague or negative phrasing before retrieval. Keeping one query and
one retrieval step makes the experiments reproducible.

\section{Reviewed-item retrieval}
\citet{rird2023} defined reviewed-item retrieval and released RIRD, the dataset used in our
English track. Items are represented only by their reviews, and the system ranks items for a
free-text query. The authors compare two strategies: combining the reviews into one item
representation first (early fusion), or scoring each review and combining the top $K$ scores
afterwards (late fusion). Late fusion works better. Their best model is a BERT encoder
fine-tuned with self-supervised contrastive learning, which reaches MAP 0.626. They also show
that an off-the-shelf dense model already beats BM25. Our English result repeats this finding
with a different, multilingual encoder, and adds two controls that the original study does not
include: identical units for both models, and a review-volume control.

\citet{korikov2024multiaspect} extend this work to queries that ask for several things at once,
such as ``a quiet place with vegan options''. They show that top-$K$ late fusion struggles when
the parts of a query appear in different reviews, and they use a large language model to split
the query into aspects that are retrieved separately and fused. We cite but do not reimplement
this approach. Our question is whether fusion is fair, not how to decompose queries, and an LLM
step would add cost and make the Thai experiments harder to reproduce. Their finding does
suggest that our Thai combination queries, such as \thai{อะไรแซ่บๆ เปรี้ยวเผ็ดเค็ม}
(something zesty: sour, spicy and salty), will be the hardest, so we report them separately.

\section{Score aggregation and the volume problem}
Turning many document scores into one entity score is an older problem from expert search,
where a person must be ranked from the documents they wrote. \citet{balog2006formal} compare
building one profile per candidate, which corresponds to the concatenated restaurant document
in our proposal, with scoring each document and aggregating the scores, which is what our final
system does. The second approach performs better. \citet{macdonald2008voting} describe these
aggregations as voting techniques and show that candidates with many documents gain an unfair
advantage unless the scores are normalised. Our 400-review cap is the reviewed-item version of
this normalisation. Finally, \citet{chen2024densex} show that the size of the retrieval unit
matters for dense retrievers, and that smaller units often help off-the-shelf models. This
motivates our comparison of review-level and sentence-level units. To our knowledge, no
previous reviewed-item retrieval study has tested sentence units or a volume control.

\section{Lexical and dense retrieval}
BM25 ranks text by term frequency, inverse document frequency and document-length
normalisation \citep{robertson2009bm25}. It is fast, deterministic and easy to interpret, and it
remains a strong baseline. Its main weakness is vocabulary mismatch: ``affordable'' and
``cheap'' express the same need but share no words.

Dense retrieval addresses this weakness by mapping queries and text into a shared vector
space. BERT showed that bidirectional Transformer pre-training produces strong contextual
representations \citep{devlin2019bert}. Sentence-BERT made sentence-level similarity practical
with Siamese encoders \citep{reimers2019sbert}, and dense passage retrieval showed that
independently encoded queries and passages can outperform BM25 for open-domain question
answering \citep{karpukhin2020dpr}. \citet{reimers2020multilingual} extended sentence encoders
to many languages through knowledge distillation. We use one of these multilingual models,
\texttt{paraphrase-multilingual-mpnet-base-v2}, so that exactly the same model can be applied to
English and Thai. This model is trained for paraphrase similarity rather than for retrieval,
which is relevant when interpreting the Thai results.

A bi-encoder encodes the query and the text separately, so it can miss the relation between
them. Negation is a clear case: ``no waiting line'' and ``a long waiting line'' are close in
embedding space. \citet{weller2024nevir} show that most neural retrievers perform close to
random on negation, and that cross-encoders, which read the query and the text together, perform
best. This motivates the selective cross-encoder reranking in our live demo.

\section{Explainable recommendation}
\citet{zhang2020explainable} survey explainable recommendation and stress that explanations
should be connected to both the user's preference and the recommendation decision. Selecting
existing review sentences as explanations is an established approach; EXTRA, for example,
frames explanation as ranking candidate review sentences \citep{li2021extra}. Extracted
sentences cannot invent facts, which is an advantage over generated explanations.

Showing a sentence does not prove that it explains the ranking. \citet{jacovi2020faithfulness}
argue that faithfulness (does the explanation reflect what the model did?) and plausibility
(does a person find it convincing?) are different properties that must be measured separately.
ERASER measures faithfulness by removing the explanation and checking how much the model output
changes \citep{deyoung2020eraser}; our deletion test applies this idea to restaurant ranking.
\citet{zhang2025enhancing} evaluate explanations from the user's point of view, including whether
the explanation is factually supported by the item; our human support rating asks a similar
question. We use both a deletion test and a human rating on the same sentences.

\section{Thai language processing and retrieval}
Thai is written without spaces between words, and sentence boundaries are not always marked by
punctuation. Sentence and word segmentation are therefore part of the retrieval pipeline. We use
PyThaiNLP \citep{phatthiyaphaibun2023pythainlp}, with the \texttt{crfcut} sentence segmenter
and the \texttt{newmm} word tokeniser. Segmentation is imperfect: \citet{lowphansirikul2020scb}
report a sentence-boundary F1 of about 0.77 for \texttt{crfcut} on product reviews, which are
informal text similar to Wongnai reviews. WangchanBERTa shows the value of
Thai-specific pre-training \citep{lowphansirikul2021wangchanberta}; we deliberately keep one
multilingual encoder so that the English and Thai experiments use the same model.

Existing Thai retrieval benchmarks cover Wikipedia, question answering or legal text rather than
reviews. On the Thai part of MIRACL \citep{zhang2023miracl}, BM25 (nDCG@10 0.484) outperforms
the multilingual dense retriever mDPR (0.358). Dense retrieval therefore cannot be assumed to win
in Thai. We found no existing Thai test set for taste-based search over restaurant reviews, which
is why we built one (Chapter~\ref{ch:data}).

\section{Positioning of this project}
\begin{table}[H]
\centering\small
\caption{Comparison of the closest related work with this project. ``Partly'' means the
property is present in a limited form.}
\label{tab:related-compare}
\begin{tabularx}{\textwidth}{@{}p{3.6cm}*{6}{>{\centering\arraybackslash}X}@{}}
\toprule
Work & Free-text query & Item ranking & Review evidence shown & Evidence evaluated & Unit and volume control & Thai\\
\midrule
Rating-based recommenders \citep{koren2009matrix,hasan2025reviewsurvey} & -- & \checkmark & Partly & -- & -- & --\\
RIRD \citep{rird2023} & \checkmark & \checkmark & -- & -- & -- & --\\
Multi-aspect RIR \citep{korikov2024multiaspect} & \checkmark & \checkmark & -- & -- & -- & --\\
EXTRA \citep{li2021extra} & -- & -- & \checkmark & Partly & -- & --\\
MIRACL Thai \citep{zhang2023miracl} & \checkmark & -- & -- & -- & -- & \checkmark\\
\textbf{This project} & \checkmark & \checkmark & \checkmark & \checkmark & \checkmark & Partly\\
\bottomrule
\end{tabularx}
\end{table}

Table~\ref{tab:related-compare} summarises the gap. Reviewed-item retrieval work ranks
restaurants for free-text queries but does not show or evaluate evidence, and does not control
for units or review volume. Explanation work ranks review sentences but not for a free-text
query over items. Thai retrieval benchmarks do not cover reviews or taste. This project brings
the three together. It replicates the RIRD dense-versus-BM25 result under a fair unit control
and a volume control; it tests whether the shown sentences are faithful and supportive; and it
applies the same retrieval step to a new human-labelled Thai taste test set. The Thai entry is
``Partly'' because the Thai evaluation is at sentence level, not restaurant level.

We do not claim to be the first to rank review sentences as explanations or to compare BM25
with dense retrieval on RIRD. The new parts are the controlled comparison (same units, volume
cap, sentence units), the combined faithfulness and support evaluation on a reviewed-item ranker,
and the Thai taste test set.
